\documentclass[12pt, a4paper]{article}

% Paquetes requeridos para el circuito
\usepackage[utf8]{inputenc}
\usepackage[T1]{fontenc}
\usepackage{circuitikz} 

\begin{document}

\begin{center}
    \begin{circuitikz}[american, scale=1.0, transform shape]
        
        % 1. Raíl inferior común (nodo 'b')
        \draw (0,0) -- (10.5,0) node[circ] {} node[right] {b};
        
        % 2. Rama V1 entre 'b' y 'c'
        \draw (0,0) node[circ] {} to[V, l=$V_1$] (0,4) node[circ] {} node[above] {c};
        
        % 3. Rama R1 horizontal entre 'c' y 'd'
        \draw (0,4) to[R, l=$R_1$] (3,4) node[circ] {} node[above] {d};
        
        % 4. Rama con fuente independiente de corriente i1 entre 'd' y 'b'
        \draw (3,4) to[I, l=$i_1$] (3,0) node[circ] {};
        
        % 5. Rama R2 horizontal entre 'd' y 'e'
        \draw (3,4) to[R, l=$R_2$] (6,4) node[circ] {} node[above] {e};
        
        % 6. Rama V2 entre 'e' y 'b'
        \draw (6,4) to[V, l=$V_2$] (6,0) node[circ] {};
        
        % 7. Etapa de salida: Fuente dependiente entre 'f' y 'b'
        \draw (8.5,4) node[circ] {} node[above] {a} to[cI, l=$\beta i_1$] (8.5,0) node[circ] {};
        
        % 8. Resistencia R3 en paralelo (entre 'f' y 'b')
        \draw (8.5,4) -- (10.5,4) to[R, l=$R_3$] (10.5,0);
        
    \end{circuitikz}
\end{center}

\end{document}